\begin{prob}
    Determine the worst case time complexity of each of the recursive algorithms below.
    In each case, state the recurrence relation describing the runtime. Solve the recurrence
    relation, either by unrolling it or showing that it is the same as a recurrence we have
    encountered in lecture.

    \begin{subprobset}
        \begin{subprob}
            \begin{minted}[autogobble]{python}
                import math
                def summation(numbers):
                    """Given a list, returns the sum of the numbers in the list."""
                    n = len(numbers)
                    if n == 0:
                        return 0
                    if n == 1:
                        return numbers[0]
                    middle = math.floor(n / 2)
                    return summation(numbers[:middle]) + summation(numbers[middle:])
            \end{minted}

            \emph{Hint:} remember that slicing causes a copy to happen!

            \begin{soln}
                Note that the slicing in \python{numbers[:middle]} takes linear time
                in the length of the array, so the recurrence relation is:
                \[
                    T(n) = \Theta(n) + 2T(n/2)
                \]
                This is the same recurrence relation as for mergesort, and the solution
                is $\Theta(n \log n)$.
            \end{soln}
        \end{subprob}


        \begin{subprob}
            \begin{minted}[autogobble]{python}
                import math
                def summation_2(numbers, start, stop):
                    """Computes the sum of numbers[start:stop]"""
                    if stop <= start:
                        return 0
                    if stop - start == 1:
                        return numbers[start]
                    middle = math.floor((start + stop) / 2)
                    left_sum = summation_2(numbers, start, middle)
                    right_sum = summation_2(numbers, middle, stop)
                    return left_sum + right_sum
            \end{minted}

            Hint: recall the formula for a geometric sum: $\sum_{p = 0}^N b^p = \frac{1 - b^{N+1}}{1-b}$

            \begin{soln}
                No slicing is being done, so a constant amount of time is required outside
                of the recursive calls. Therefore the recurrence relation is:
                \[
                    T(n) = 2T(n/2) + \Theta(1)
                \]

                We have not seen this recurrence before, so we must solve it by unrolling:
                \begin{align*}
                    T(n)
                    &=
                        2T(n/2) + 1
                        \\
                    &=
                        2\left[
                            2T(n/4) + 1
                        \right]
                        +1
                        \\
                    &=
                        4 T(n/4) + 2 + 1
                        \\
                    &=
                        4 \left[
                            2 T(n/8) + 1
                        \right]
                        + 2 + 1
                        \\
                    &=
                        8 T(n/8) + 4 + 2 + 1
                        \\
                \end{align*}
                At this point, a pattern has emerged, and we can guess that on unroll
                $k$ we'll have:
                \[
                    T(n) = 2^k \cdot T(n/2^k) + 1 + 2 + 4 + \ldots + 2^{k-1}
                \]
                We will reach $T(1)$ when $k = \log_2 n$. Making this substitution, we have:
                \begin{align*}
                    T(n) &= 2^{\log_2 n} \cdot T(1) + 1 + 2 + 4 + \ldots + 2^{(\log_2 n) - 1}\\
                         &= n + \left(1 + 2 + 4 + \ldots + 2^{(\log_2 n) - 1}\right)
                 \intertext{The summation is a geometric sum, and, using the hint with $b = 2$, we have:}
                         &= n + \frac{1 - 2^{\log_2 n}}{1 - 2}
                            \\
                         &= n + \frac{1 - n}{1 - 2}
                            \\
                         &=
                            \Theta(n)
                \end{align*}
            \end{soln}
        \end{subprob}

        \begin{subprob}
            In this subproblem, you may assume that \python{arr} is sorted.
            \begin{minted}[autogobble]{python}
                import math
                def ternary_search(arr, t, start, stop):
                    """Return the index of t in arr[start:stop]"""
                    if stop - start <= 0:
                        return None
                    n = stop - start
                    left_ix = math.floor(start + (stop - start) / 3)
                    right_ix = math.floor(start + 2 * (stop - start) / 3)
                    if arr[left_ix] == t:
                        return left_ix
                    if arr[right_ix] == t:
                        return right_ix
                    if t < arr[left_ix]:
                        return ternary_search(arr, t, start, left_ix)
                    elif t > arr[right_ix]:
                        return ternary_search(arr, t, right_ix + 1, stop)
                    else:
                        return ternary_search(arr, t, left_ix + 1, right_ix)
            \end{minted}

            \begin{soln}
                This is like binary search, but the recursive calls are on arrays of
                size $n/3$. Hence the recurrence is:
                \[
                    T(n) = T(n/3) + \Theta(1)
                \]
                The solution will be the same as the solution for binary search, except
                all instances of $\log_2 n$ will be replaced by $\log_3 n$, so the overall
                time complexity is the same: $\Theta(\log n)$.
            \end{soln}
        \end{subprob}

    \end{subprobset}

\end{prob}
